\documentclass[12pt,a4paper]{article}
\usepackage[UTF8]{ctex}          % 中文支持
\usepackage{verbatim}
\usepackage{fancyhdr}
\usepackage{graphicx}
\usepackage{amsfonts}
\usepackage{booktabs}
\usepackage{amsmath,amssymb}     % 数学符号
\usepackage{listings}            % 代码排版
\usepackage{xcolor}              % 颜色
\usepackage{geometry}            % 页面边距
\usepackage{hyperref}            % 超链接
\geometry{left=1.5cm,right=1.5cm,top=2.5cm,bottom=2.5cm}

% MATLAB 代码样式
\lstset{
    language=Matlab,
    basicstyle=\ttfamily\small,
    keywordstyle=\color{blue},
    commentstyle=\color{gray},
    stringstyle=\color{red},
    numbers=left,
    numberstyle=\tiny\color{gray},
    stepnumber=1,
    numbersep=5pt,
    backgroundcolor=\color{white},
    showspaces=false,
    showstringspaces=false,
    frame=single,
    tabsize=4,
    captionpos=b,
    breaklines=true,
    breakatwhitespace=false,
    escapeinside={\%*}{*)},
    morekeywords={u,v,w,RichardsonExtrapolation,romberg,sinx_over_x,stable_b,exp_m1_over_t,print_romberg}
}

\title{第九周上机作业}
\author{PB23000270 李佩优}
\date{\today}

\begin{document}

\maketitle

\section{常微分方程初值问题的数值解法}

\subsection{程序概述}

程序 \texttt{hw9.m} 实现了三种经典的数值方法：显式欧拉方法（Forward Euler）、隐式欧拉方法（Backward Euler）和梯形公式方法（Crank–Nicolson），用于求解如下常微分方程组的初值问题：
\[
y' = Ay, \quad y(0) = y_0,
\]
其中
\[
A = \begin{pmatrix} -8003 & 1999 \\ 23988 & -6004 \end{pmatrix}, \qquad
y_0 = \begin{pmatrix} 1 \\ 4 \end{pmatrix}.
\]
计算时取时间步长 \(h = 0.004\)，共进行 \(N = 12\) 步，即计算 \(y_0, y_1, \dots, y_{12}\) 共 13 个节点的数值解。程序依次执行三种方法，并将每一时间层的解以矩阵形式输出。

\subsection{算法原理}

考虑一阶常微分方程组初值问题
\[
y' = f(t, y), \quad y(t_0) = y_0.
\]
对于线性自治系统 \(y' = Ay\)，右端项 \(f(t,y) = Ay\) 不显含时间 \(t\)，三种单步法的迭代格式如下：

\subsubsection{显式欧拉方法}
\[
y_{n+1} = y_n + h A y_n = (I + hA) y_n,
\]
其中 \(I\) 为与 \(A\) 同阶的单位矩阵。该方法为显式格式，每步仅需一次矩阵-向量乘法，计算量小，但稳定性条件较苛刻。

\subsubsection{隐式欧拉方法}
\[
y_{n+1} = y_n + h A y_{n+1} \quad \Longrightarrow \quad (I - hA) y_{n+1} = y_n.
\]
每一步需要求解一个线性方程组，对二维问题可直接求逆或使用反斜杠算子，计算成本略高于显式方法，但具有更好的稳定性。

\subsubsection{梯形公式方法（Crank–Nicolson）}
\[
y_{n+1} = y_n + \frac{h}{2} \bigl(A y_n + A y_{n+1}\bigr) \quad \Longrightarrow \quad \bigl(I - \tfrac{h}{2}A\bigr) y_{n+1} = \bigl(I + \tfrac{h}{2}A\bigr) y_n.
\]
该方法在时间方向具有二阶精度，且对于线性问题是无条件稳定的（若 \(A\) 的特征值实部非正），常优于前两种一阶方法。

\subsection{程序结构说明}

程序的主流程：首先定义矩阵 \(A\)、初值 \(y_0\)、步长 \(h\) 和步数 \(N\)，然后依次调用三种方法进行计算并显示结果。
为便于矩阵运算，预先计算出所需的系数矩阵，再通过循环递推生成解序列。
在北太天元和matlab上均运行良好。核心代码结构如下：

\begin{lstlisting}[caption={主程序框架}, label={lst:hw9main}]
clear; clc; close all;
format short e

A = [-8003, 1999; 23988, -6004] ;
y0 = [1; 4] ;
h = 0.004 ;
N = 12 ;

% 显式欧拉方法
I_plus_A = eye(2) + h * A;
M = zeros(2, N+1);
M(:, 1) = y0;
for i = 1:N
    y = I_plus_A * M(:, i) ;
    M(:, i+1) = y;
end
fprintf('显式欧拉方法：\n')
disp(M)

% 隐式欧拉方法
I_minus_A = eye(2) - h * A;
M = zeros(2, N+1);
M(:, 1) = y0;
for i = 1:N
    y = I_minus_A \ M(:, i);
    M(:, i+1) = y;
end
fprintf('隐式欧拉方法：\n')
disp(M)

% 梯形公式方法
I_minus_A = eye(2) - (h/2) * A;
I_plus_A = eye(2) + (h/2) * A;
M = zeros(2, N+1);
M(:, 1) = y0;
for i = 1:N
    y = I_minus_A \ (I_plus_A * M(:, i));
    M(:, i+1) = y;
end
fprintf('梯形公式方法：\n')
disp(M)
\end{lstlisting}

三种方法均使用 \texttt{eye(2)} 生成 \(2\times2\) 单位矩阵，并预计算迭代所需的矩阵，使得循环体内仅进行矩阵-向量乘法或线性方程组求解，效率高且代码简洁。输出矩阵 \texttt{M} 的每一列依次对应 \(y_0, y_1, \dots, y_{12}\)。

\subsection{数值结果与分析}

为直观对比，将三种方法的数值解分别列于表\ref{tab:forward}、表\ref{tab:backward}和表\ref{tab:cn}（以普通小数形式给出，不使用科学记数法）。每个表格中，第一行为时间步 \(n = 0,1,\dots,12\)，第二、三行分别对应向量 \(y_n\) 的两个分量。

\begin{table}[htbp]
\centering
\caption{显式欧拉方法数值解}
\label{tab:forward}
\resizebox{\textwidth}{!}{%
\begin{tabular}{c|ccccccccccccc}
\toprule
$n$ & 0 & 1 & 2 & 3 & 4 & 5 & 6 & 7 & 8 & 9 & 10 & 11 & 12 \\
\midrule
$y_1$ & 1.0000 & 0.97200 & 0.94478 & 0.91833 & 0.89262 & 0.86762 & 0.84333 & 0.81977 & 0.79373 & 0.94141 & -8.4297 & 505.77 & -27776 \\
$y_2$ & 4.0000 & 3.8880 & 3.7791 & 3.6733 & 3.5705 & 3.4705 & 3.3733 & 3.2787 & 3.1962 & 2.5970 & 30.559 & -1512.2 & 83334 \\
\bottomrule
\end{tabular}
}
\end{table}

\begin{table}[htbp]
\centering
\caption{隐式欧拉方法数值解}
\label{tab:backward}
\resizebox{\textwidth}{!}{%
\begin{tabular}{c|ccccccccccccc}
\toprule
$n$ & 0 & 1 & 2 & 3 & 4 & 5 & 6 & 7 & 8 & 9 & 10 & 11 & 12 \\
\midrule
$y_1$ & 1.0000 & 0.97276 & 0.94627 & 0.92049 & 0.89542 & 0.87103 & 0.84731 & 0.82423 & 0.80178 & 0.77994 & 0.75870 & 0.73803 & 0.71793 \\
$y_2$ & 4.0000 & 3.8911 & 3.7851 & 3.6820 & 3.5817 & 3.4841 & 3.3892 & 3.2969 & 3.2071 & 3.1198 & 3.0348 & 2.9521 & 2.8717 \\
\bottomrule
\end{tabular}
}
\end{table}

\begin{table}[htbp]
\centering
\caption{梯形公式方法数值解}
\label{tab:cn}
\resizebox{\textwidth}{!}{%
\begin{tabular}{c|ccccccccccccc}
\toprule
$n$ & 0 & 1 & 2 & 3 & 4 & 5 & 6 & 7 & 8 & 9 & 10 & 11 & 12 \\
\midrule
$y_1$ & 1.0000 & 0.97239 & 0.94554 & 0.91943 & 0.89404 & 0.86935 & 0.84534 & 0.82200 & 0.79930 & 0.77723 & 0.75577 & 0.73490 & 0.71461 \\
$y_2$ & 4.0000 & 3.8895 & 3.7821 & 3.6777 & 3.5762 & 3.4774 & 3.3814 & 3.2880 & 3.1972 & 3.1089 & 3.0231 & 2.9396 & 2.8584 \\
\bottomrule
\end{tabular}
}
\end{table}

\textbf{分析}：

矩阵 \(A\) 的特征值可通过计算得到，约为 \(\lambda_1 \approx -5.003,\ \lambda_2 \approx -14004\)。两个特征值均为负实数，且模值相差悬殊，表明该系统是一个刚性（stiff）问题。显式欧拉方法的绝对稳定区间为复平面上的单位圆 \(|1 + h\lambda| < 1\)。对于 \(\lambda_2 \approx -14004\)，取 \(h = 0.004\) 得到 \(h\lambda_2 \approx -56.016\)，远超出稳定域，因此数值解迅速发散：由表\ref{tab:forward}可见，从第9步开始解分量出现剧烈振荡并迅速爆炸至 \(10^4\) 量级，完全失去物理意义。

隐式欧拉方法和梯形公式方法均为 A-稳定的（梯形公式是 A-稳定的，隐式欧拉是 L-稳定的），对于任意 \(h>0\) 只要 \(\operatorname{Re}(\lambda) \le 0\) 数值解均能保持稳定。从表\ref{tab:backward}和表\ref{tab:cn}可见，隐式欧拉和梯形公式的解均随时间单调衰减，两个分量平稳地从初值 \((1,4)^\top\) 向零方向衰减，且未出现任何振荡或发散。梯形公式作为二阶方法，在每一步的精度上略优于一阶的隐式欧拉，但由于步长较大、系统刚性极强，两种方法的解曲线在多数节点上差异不大，均能正确反映解的定性行为。

\subsection{总结}

本实验针对一个刚性线性常微分方程组，分别采用显式欧拉、隐式欧拉和梯形公式进行数值求解。实验结果显示，对于所选的步长 \(h=0.004\)，显式欧拉方法因步长超出稳定性限制而发散，而隐式欧拉和梯形公式方法凭借其优越的稳定性，给出了合理的数值解。该对比直观地展示了数值方法稳定性在刚性方程求解中的关键作用：单纯的精度考虑不足以保证计算成功，必须选取与问题刚性相匹配的稳定方法（如 A-稳定的隐式或梯形公式）。此外，程序结构简洁、高效，利用 MATLAB 矩阵运算的特性使代码清晰易读，为更复杂的常微分方程组数值求解提供了良好的实现模板。

\newpage
\section{计算机代码}

\small
\hrule
\verbatiminput{hw9.m}
\hrule

\newpage
\section{原始输出}
以下为 MATLAB 命令窗口中直接输出的原始结果（使用 \texttt{format short e}）。

\begin{verbatim}
显式欧拉方法：
  列 1 至 7

   1.0000e+00   9.7200e-01   9.4478e-01   9.1833e-01   8.9262e-01   8.6762e-01   8.4333e-01
   4.0000e+00   3.8880e+00   3.7791e+00   3.6733e+00   3.5705e+00   3.4705e+00   3.3733e+00

  列 8 至 13

   8.1977e-01   7.9373e-01   9.4141e-01  -8.4297e+00   5.0577e+02  -2.7776e+04
   3.2787e+00   3.1962e+00   2.5970e+00   3.0559e+01  -1.5122e+03   8.3334e+04

隐式欧拉方法：
  列 1 至 7

   1.0000e+00   9.7276e-01   9.4627e-01   9.2049e-01   8.9542e-01   8.7103e-01   8.4731e-01
   4.0000e+00   3.8911e+00   3.7851e+00   3.6820e+00   3.5817e+00   3.4841e+00   3.3892e+00

  列 8 至 13

   8.2423e-01   8.0178e-01   7.7994e-01   7.5870e-01   7.3803e-01   7.1793e-01
   3.2969e+00   3.2071e+00   3.1198e+00   3.0348e+00   2.9521e+00   2.8717e+00

梯形公式方法：
  列 1 至 7

   1.0000e+00   9.7239e-01   9.4554e-01   9.1943e-01   8.9404e-01   8.6935e-01   8.4534e-01
   4.0000e+00   3.8895e+00   3.7821e+00   3.6777e+00   3.5762e+00   3.4774e+00   3.3814e+00

  列 8 至 13

   8.2200e-01   7.9930e-01   7.7723e-01   7.5577e-01   7.3490e-01   7.1461e-01
   3.2880e+00   3.1972e+00   3.1089e+00   3.0231e+00   2.9396e+00   2.8584e+00
\end{verbatim}

\end{document}